\label{sec:order-distance-length}

\begin{topicbox}{Distance on the Line}
\begin{exposition}
The absolute value of a difference turns $\mathbb{R}$ into a metric space. Length of an interval is the distance between its endpoints. These are the quantitative tools underneath every later definition of nearness.
\end{exposition}
\end{topicbox}
\begin{definitionbox}{Definition (Distance on the Real Line)}
\begin{definition}[Distance on the Real Line]
\label{def:distance-on-real-line}
The \emph{distance} on $\mathbb{R}$ is the map
\[
  d:\mathbb{R}\times\mathbb{R}\longrightarrow\mathbb{R}
\]
defined by
\[
  d(x,y)=|x-y|
  \qquad(x,y\in\mathbb{R}).
\]
\end{definition}
\end{definitionbox}
\begin{remark*}[Standard quantified statement]
\[
\forall x,y\in\mathbb{R}\;(d(x,y)=|x-y|).
\]
\end{remark*}
\begin{remark*}[Predicate reading]
\[
\operatorname{Distance}(d,\mathbb{R})
\Longleftrightarrow
\forall x,y\in\mathbb{R}\;(d(x,y)=|x-y|).
\]
\end{remark*}
\begin{remark*}[Interpretation]
The distance between two points is the length of the segment joining them. It is the metric that makes ``close'' precise and is the basis for balls, neighborhoods, and open sets.
\end{remark*}
\begin{dependencies}
\begin{itemize}
  \item \hyperref[thm:triangle-inequality]{Triangle Inequality}
\end{itemize}
\end{dependencies}
\begin{definitionbox}{Definition (Length of a Bounded Interval)}
\begin{definition}[Length of a Bounded Interval]
\label{def:length-of-interval}
Let $a,b\in\mathbb{R}$ with $a\leq b$.  The \emph{length} of any bounded
interval with endpoints $a$ and $b$ is
\[
  \ell=b-a=d(a,b).
\]
\end{definition}
\end{definitionbox}
\begin{remark*}[Standard quantified statement]
\[
\forall a,b\in\mathbb{R}\;\bigl(a\leq b \Longrightarrow \ell([a,b])=b-a=d(a,b)\bigr).
\]
\end{remark*}
\begin{remark*}[Predicate reading]
\[
\operatorname{Length}([a,b],\,b-a)
\Longleftrightarrow
a\leq b.
\]
\end{remark*}
\begin{remark*}[Interpretation]
Length depends only on the endpoints, not on whether they are included. It is the distance $d(a,b)$ and measures the size of the basic pieces of the line.
\end{remark*}
\begin{dependencies}
\begin{itemize}
  \item \hyperref[def:distance-on-real-line]{Distance on the Real Line}
\end{itemize}
\end{dependencies}
\begin{theorembox}{Theorem (The Distance Function Is a Metric)}
\begin{theorem}[The Distance Function Is a Metric]
\label{thm:distance-is-a-metric}
Let $d:\mathbb{R}\times\mathbb{R}\longrightarrow\mathbb{R}$ be defined by
$d(x,y)=|x-y|$.  Then, for all $x,y,z\in\mathbb{R}$,
\[
  d(x,y)\geq0,
  \qquad
  d(x,y)=0\Longleftrightarrow x=y,
\]
\[
  d(x,y)=d(y,x),
\]
and
\[
  d(x,z)\leq d(x,y)+d(y,z).
\]
Hence $(\mathbb{R},d)$ is a metric space.
\hyperref[prf:distance-is-a-metric]{Go to proof.}
\end{theorem}
\end{theorembox}
\begin{remark*}[Standard quantified statement]
\[
\forall x,y,z\in\mathbb{R}\;\bigl(d(x,y)\geq 0\ \land\ (d(x,y)=0\Longleftrightarrow x=y)\ \land\ d(x,y)=d(y,x)\ \land\ d(x,z)\leq d(x,y)+d(y,z)\bigr).
\]
\end{remark*}
\begin{remark*}[Predicate reading]
\[
d:\mathbb{R}\times\mathbb{R}\longrightarrow\mathbb{R},\qquad
\mathsf{MetricSpace}(\mathbb{R},d)
\Longleftrightarrow
\text{$d$ satisfies nonnegativity, identity of indiscernibles, symmetry, and the triangle inequality.}
\]
\end{remark*}
\begin{remark*}[Interpretation]
The three metric axioms are exactly the properties of $|\cdot|$: nonnegativity with the zero-distance characterization, symmetry, and the triangle inequality. This is the precise sense in which $\mathbb{R}$ is a one-dimensional metric space.
\end{remark*}
\begin{dependencies}
\begin{itemize}
  \item \hyperref[def:distance-on-real-line]{Distance on the Real Line}
  \item \hyperref[thm:triangle-inequality]{Triangle Inequality}
\end{itemize}
\end{dependencies}

\begin{theorembox}{Theorem (Real Line Structural Order Facts)}
\begin{theorem}[Real Line Structural Order Facts]\label{thm:real-line-structural-order-facts}
The real line has the following structural properties: $\mathbb{R}$ is
linearly ordered by $\leq$, $\mathbb{R}$ is an ordered field, and
$\mathbb{R}$ satisfies the least upper bound property.  Explicitly, whenever
$A\subseteq\mathbb{R}$ is nonempty and bounded above in $\mathbb{R}$, there
exists $s\in\mathbb{R}$ such that
\[
  (\forall a\in A)(a\leq s)
\]
and
\[
  (\forall u\in\mathbb{R})
  \left[
  (\forall a\in A)(a\leq u)
  \Longrightarrow
  s\leq u
  \right].
\]
\hyperref[prf:real-line-structural-order-facts]{Go to proof.}
\end{theorem}
\end{theorembox}

\begin{remark*}[Standard quantified statement]
\[
\operatorname{TotallyOrderedSet}(\mathbb{R},\leq)
\land
\operatorname{OrderedField}(\mathbb{R})
\land
\operatorname{HasLeastUpperBoundProperty}(\mathbb{R}).
\]
\end{remark*}

\begin{remark*}[Interpretation]
The real line combines order, algebra, and completeness. These three
structures are the ambient facts used by distance, interval, and convergence
arguments throughout the volume.
\end{remark*}

\begin{dependencies}
\begin{itemize}
  \item \hyperref[thm:r-totally-ordered-set]{$\mathbb{R}$ Is a Totally Ordered Set}
  \item \hyperref[thm:r-is-an-ordered-field]{$\mathbb{R}$ Is an Ordered Field}
  \item \hyperref[thm:lub-property-r]{Least-Upper-Bound Property}
\end{itemize}
\end{dependencies}
